\documentclass[11pt]{article}
\usepackage[margin=1in]{geometry}
\usepackage{booktabs}
\usepackage{array}
\usepackage{xcolor}
\usepackage{hyperref}
\usepackage{enumitem}
\usepackage{parskip}

\hypersetup{
  colorlinks=true,
  linkcolor=blue!50!black,
  urlcolor=blue!50!black,
}

\newcommand{\flag}[1]{\textcolor{red!70!black}{\textbf{#1}}}
\newcommand{\fact}[0]{\textcolor{green!40!black}{[fact]}}
\newcommand{\guess}[0]{\textcolor{orange!80!black}{[hypothesis]}}

\title{Colombia Validation: Pipeline vs.\ Samuel's Published Results\\
\large World Bank climate-health burden projections}
\author{Aaron Osgood-Zimmerman}
\date{2026-05-07}

\begin{document}
\maketitle

\section*{TL;DR}

Running the global R pipeline in \texttt{COLOMBIA\_VERIFICATION} mode against
Samuel Cuervo's published Colombia 2010--2019 results, we now pass
\textbf{18 of 19} headline metrics at $\pm 5\%$ tolerance. National totals
agree to $-2.26\%$ (non-optimal deaths), $-2.35\%$ (heat), $-2.24\%$ (cold),
and $-3.35\%$ (non-optimal years of life lost, YLL). Top-3 cause rankings
match exactly for both heat and cold across deaths, YLLs, and population
attributable fractions (PAFs). The only remaining failure is the heat YLL
for road injuries (\texttt{inj\_trans\_road}) at $-11.76\%$. Section
\ref{sec:hypotheses} describes hypotheses for the residual gap; Section
\ref{sec:verif-vs-prod} describes what is stripped out before the pipeline
runs in production / global mode.

This document is a snapshot as of 2026-05-07 of the
\texttt{verification-depto-day} branch.

\section{Background and terminology}

\paragraph{Project.} World Bank statistical consulting work quantifying
mortality and morbidity attributable to non-optimal ambient temperature.
Validating against Colombia 2010--2019 (location\_id 125) before scaling
globally and projecting to 2100 under SSP (Shared Socioeconomic Pathway)
scenarios. Built on the Burkart et al.\ (\textit{Lancet} 2021) methodology
and the GBD (Global Burden of Disease) 2023 framework.

\paragraph{Reference.} Samuel Cuervo's Colombia work used a 12-script R
pipeline plus pre-processed input data (DANE vital registration, ERA5
temperature, WorldPop). His published Colombia totals for 2010--2019 are
the targets of this validation.

\paragraph{Key terms.} Defined inline at first use:
\begin{description}[style=multiline,leftmargin=3.5cm]
  \item[ERF] Exposure-Response Function. The relative-risk (RR) curve mapping
    daily temperature to cause-specific RR for one of 17 temperature-sensitive
    causes. Provided as 1000 posterior draws per (zone, daily temperature, cause).
  \item[TMREL] Theoretical Minimum Risk Exposure Level. The death-weighted
    optimal temperature; days hotter than the TMREL count as ``heat,''
    cooler as ``cold.''
  \item[PAF] Population Attributable Fraction. Fraction of deaths in a
    population that would not have occurred under counterfactual exposure
    (here, all days at the TMREL).
  \item[SEV] Summary Exposure Value. A standardized $[0,1]$ exposure metric.
    In Burkart, diagnostic only; in Samuel's Colombia code (and our
    verification mode) it enters the burden formula multiplicatively.
  \item[YLL] Years of Life Lost. Attributable deaths $\times$ remaining
    life expectancy at age of death.
  \item[depto] Departamento --- Colombia's admin-1 / state-level
    subdivision. Carried as \texttt{subloc\_id} (zero-padded
    DIVIPOLA codes, e.g.\ ``05'' for Antioquia).
\end{description}

\section{Validation methodology}

\subsection{Verification mode toggles}
\label{subsec:toggles}

The pipeline has a \texttt{COLOMBIA\_VERIFICATION} flag (\texttt{config.R}) that,
when \texttt{TRUE}, replicates six methodological choices specific to Samuel's
Colombia code. A safety guard refuses to run with this flag set if
\texttt{LOCATION\_ID != 125}.

\begin{enumerate}
  \item \textbf{No RR rescaling at the TMREL.} Burkart rescales each ERF curve
    so $RR(\text{TMREL}) = 1$ before computing PAFs; Samuel uses the raw
    posterior-mean RR. Verification mode skips the rescale.
  \item \textbf{Per-row PAF floor at zero.} Samuel zeroes out negative
    contributions to the PAF before summing within (year, depto, cause, risk).
    Burkart sums signed contributions and lets cancellation happen.
  \item \textbf{TMREL averaged across years.} TMREL files contain values for
    1990, 2010, and 2020. Verification mode averages all three with equal
    weight to produce a single per-zone value used for all study years.
    Production mode uses the year-specific TMREL with nearest-year fill for
    missing years.
  \item \textbf{Burden formula multiplied by SEV.} Verification mode computes
    attributable burden as $\text{deaths} \times \text{PAF} \times \text{SEV}$
    rather than $\text{deaths} \times \text{PAF}$.
  \item \textbf{Samuel's daily-summation SEV bug.} Samuel's
    \texttt{11\_carga\_atribuible.R:444--469} sums daily SEV contributions
    over 365 days and then caps at 1, which inflates the SEV by roughly
    a factor of $N_{\text{days}}/\text{year}$ before truncation. Verification
    mode replicates this bug exactly.
  \item \textbf{Department-specific life tables; \texttt{ex} adjustments.}
    YLL computation uses department-specific (not national) DANE life tables;
    age 0 is remapped to age 1 for the 0--4 group; remaining life expectancy
    is reduced by 2.5 years for ages $<80$ and held at a fixed 10 years for
    ages $\geq 80$, mirroring Samuel's
    \texttt{11\_carga\_atribuible.R:570--577}.
\end{enumerate}

\subsection{Subnational refactor}

The pipeline computes PAFs at \texttt{(year, subloc\_id, cause, risk)}
granularity and merges with mortality at full
\texttt{(year, subloc\_id, age\_group\_id, sex\_id, cause)} granularity.
Pre-refactor it computed PAFs only at \texttt{(year, cause)} and merged
against country-level mortality, losing the within-country covariance
between high-mortality and high-temperature departments.

\subsection{Depto-day verification branch}
\label{subsec:depto-day}

The most recent change is a daily-attribution branch
(\texttt{05\_compute\_pafs.R}, gated on \texttt{COLOMBIA\_VERIFICATION}):

\begin{enumerate}
  \item Aggregate temperature to (date, depto, zone, daily-temp bin) and
    compute within-(date, depto) population fractions $\pi$.
  \item Compute daily PAFs as $\sum_i \max(0, \pi_i (RR_i-1)/RR_i)$ within
    (date, depto, cause, risk).
  \item Merge daily PAFs with daily mortality at (date, depto, cause),
    broadcasting across age and sex.
  \item Multiply by the annual SEV at (year, depto, cause).
  \item Sum days $\to$ annual burden at
    (year, depto, age, sex, cause).
\end{enumerate}

This captures the within-year covariance between hot days and high-mortality
days that Samuel's daily-attribution pattern picks up but the
annual-PAF $\times$ annual-deaths formula does not.

\subsection{Comparison layers}

\texttt{util\_compare\_to\_samuel.R} compares 19 metrics in five layers:
national totals, per-million rates, YLL totals, sex split, and per-cause
YLLs. Top-3 cause rankings are also checked. Per-cause PAFs are reported
but excluded from the pass/fail count because of a \fact\ metric mismatch:
the pipeline reports death-weighted PAFs, while Samuel publishes plain
pixel-day-mean PAFs --- the two are not reconcilable from his published
output alone.

\section{Numerical results}
\label{sec:results}

Table \ref{tab:results} reports all 19 headline metrics. Tolerance is $\pm 5\%$.

\begin{table}[h]
\centering
\small
\caption{Pipeline vs.\ Samuel's published Colombia 2010--2019 results,
\texttt{COLOMBIA\_VERIFICATION = TRUE}, summary mode (point-estimate RR), depto-day attribution. Tolerance $\pm 5\%$.
Source: \texttt{output/results/comparison\_to\_samuel\_125.csv} as of 2026-05-07.}
\label{tab:results}
\begin{tabular}{l r r r c}
\toprule
Metric & Pipeline & Samuel & \% diff & Verdict \\
\midrule
\multicolumn{5}{l}{\textit{National totals}} \\
Total attributable deaths (non-optimal) & 9{,}257.5 & 9{,}472 & $-2.26$ & PASS \\
Heat-attributable deaths & 2{,}247.7 & 2{,}301.7 & $-2.35$ & PASS \\
Cold-attributable deaths & 7{,}009.9 & 7{,}170.3 & $-2.24$ & PASS \\
Cold share of attributable deaths & 0.7572 & 0.7570 & $+0.03$ & PASS \\
17-cause mortality denominator & 882{,}545 & 884{,}628 & $-0.24$ & PASS \\
Share of 17-cause mortality (PAF) & 0.0105 & 0.0105 & $-0.10$ & PASS \\
\midrule
\multicolumn{5}{l}{\textit{Per-million rates (implied population $\sim$46.4M)}} \\
Rate per 1M (non-optimal) & 19.95 & 20.41 & $-2.26$ & PASS \\
Rate per 1M (heat) & 4.84 & 5.00 & $-3.14$ & PASS \\
Rate per 1M (cold) & 15.10 & 15.41 & $-1.98$ & PASS \\
\midrule
\multicolumn{5}{l}{\textit{YLL totals}} \\
Total YLL (non-optimal) & 167{,}074 & 172{,}870 & $-3.35$ & PASS \\
Heat YLL & 60{,}595 & 63{,}556 & $-4.66$ & PASS \\
Cold YLL & 106{,}479 & 109{,}313 & $-2.59$ & PASS \\
\midrule
\multicolumn{5}{l}{\textit{Sex split (deaths)}} \\
Men & 5{,}362.7 & 5{,}510 & $-2.67$ & PASS \\
Women & 3{,}894.8 & 3{,}962 & $-1.70$ & PASS \\
\midrule
\multicolumn{5}{l}{\textit{Per-cause YLL (causes Samuel reports numerically)}} \\
YLL heat --- inj\_homicide & 29{,}277 & 30{,}045 & $-2.56$ & PASS \\
YLL heat --- inj\_trans\_road & 9{,}651 & 10{,}937 & $\mathbf{-11.76}$ & \flag{FAIL} \\
YLL heat --- cvd\_ihd & 9{,}369 & 9{,}408 & $-0.41$ & PASS \\
YLL cold --- cvd\_ihd & 48{,}512 & 48{,}907 & $-0.81$ & PASS \\
YLL cold --- cvd\_stroke & 10{,}293 & 10{,}372 & $-0.76$ & PASS \\
\midrule
\multicolumn{5}{l}{\textbf{Total: 18 PASS / 1 FAIL of 19}} \\
\bottomrule
\end{tabular}
\end{table}

\paragraph{Top-3 cause rankings.} Where Samuel publishes ordered top-3 lists,
the pipeline matches exactly:
\begin{itemize}[itemsep=2pt]
  \item Heat top-3 YLL: \texttt{inj\_homicide, inj\_trans\_road, cvd\_ihd} (3/3 overlap)
  \item Cold top-3 YLL: \texttt{cvd\_ihd, resp\_copd, cvd\_stroke} (3/3 overlap)
  \item Heat top-3 PAF: \texttt{inj\_homicide, inj\_trans\_road, diabetes} (3/3 overlap)
  \item Cold top-3 PAF: \texttt{inj\_drowning, resp\_copd, cvd\_ihd} (2/3 overlap;
    Samuel has \texttt{lri} in third)
\end{itemize}

\section{Hypotheses for the residual gap}
\label{sec:hypotheses}

These are \guess{} ranked by what we would investigate first if pursuing
exact replication. Items are not numerically attributed because we have not
isolated their individual contributions.

\subsection*{H1. Mortality input is imputed upstream}

Samuel's input file is \texttt{mortalidad\_diaria\_DANE\_2010\_2019\_imput.rds};
the \texttt{imput} suffix suggests upstream missingness imputation that we
do not replicate. Our 17-cause denominator is 882{,}545 deaths vs.\ Samuel's
884{,}628 ($-0.24\%$). The EXTRANJERO filter (\texttt{codptore == 75}, deaths
of residents abroad) accounts for 1{,}987 of the 2{,}083 missing rows;
\textasciitilde 96 deaths remain unexplained. The total-level effect is small,
but if imputation is non-uniform across causes it could account for
cause-specific gaps including \texttt{inj\_trans\_road}.
\textbf{Verification path:} request the pre-imputation file or the
imputation script from Samuel.

\subsection*{H2. ERF point-estimate convention}

\texttt{01\_load\_erf.R} computes \texttt{rr\_mean} as the arithmetic mean
across 1000 posterior draws. If Samuel uses the median draw or the
publication's point estimate, the two differ when the draw distribution is
skewed --- which is plausible for the heat tail of the road-injury ERF.
\textbf{Verification path:} swap \texttt{mean()} for \texttt{median()} in the
ERF aggregation, rerun, see if \texttt{inj\_trans\_road} closes.

\subsection*{H3. TMREL provenance}

In verification mode we average IHME's 1990 / 2010 / 2020 TMREL files with
equal weight. Samuel's TMREL provenance is not fully documented in what he
sent; if his TMRELs come from a different source or were derived rather than
averaged, sub-degree shifts at the boundary push borderline days from
``optimal'' to ``heat'' or ``cold,'' which propagates into all per-cause
heat/cold splits. \textbf{Verification path:} compare our derived TMREL
file to any TMREL file in Samuel's repo we have not yet diffed.

\subsection*{H4. Garbage-code redistribution upstream of the cause map}

Our cause map is a 1:1 string mapping from Samuel's Spanish
\texttt{c\_muerte} labels to GBD acause codes. GBD typically redistributes
``garbage'' codes (R-codes / ill-defined causes) proportionally to specific
causes. If Samuel's \texttt{c\_muerte} field already has garbage-code
redistribution applied with a specific coefficient set, and we operate on
his file as-is, the totals stay close (they do: $-0.24\%$ on the denominator)
but per-cause numbers drift, which is the pattern we observe. \textbf{Verification path:} ask Samuel whether his cause field
incorporates garbage-code redistribution and which coefficients.

\subsection*{H5. Cause-specific age restriction (\texttt{inj\_trans\_road} only)}

Samuel may zero out road injuries for the 0--4 age group before
attribution; we do not. Combined with the YLL adjustment (\texttt{ex - 2.5}
for ages $<80$), even a small misassignment to 0--4 leverages heavily
because the \texttt{ex\_adj} for that group is roughly 70 years. This is
specifically the cause where we still fail.
\textbf{Verification path:} inspect his
\texttt{11\_carga\_atribuible.R} for any age filtering on
\texttt{relacion\_transporte} / \texttt{accidentes\_trafico}, or check whether his
input mortality has zero road-injury deaths at 0--4.

\subsection*{H6. Daily-mortality NA / boundary handling}

Our daily save filter is \texttt{!is.na(date) \& !is.na(acause)}. If Samuel
keeps NA-date deaths under a default-date convention rather than dropping
them, daily attribution shifts slightly. Same family of issue as H1, but
isolable on our side.

\subsection*{Comparison-script artifact (not a real gap)}

The per-cause PAF rows (excluded from the 18-of-19 count) report
death-weighted PAFs from our pipeline against plain pixel-day-mean PAFs from
Samuel. The two are not reconcilable from his published outputs alone, so
we report ours and document the metric mismatch in
\texttt{util\_compare\_to\_samuel.R}.

\section{Verification mode vs.\ production / global pipeline}
\label{sec:verif-vs-prod}

The validation just described is run with \texttt{COLOMBIA\_VERIFICATION = TRUE}.
Production / global runs set this flag to \texttt{FALSE}, which reverts to
the Burkart-correct methodology.

\subsection{What is reverted in production}

Each verification toggle from Section \ref{subsec:toggles} is undone:

\begin{enumerate}
  \item \textbf{RR is rescaled to 1.0 at the TMREL} per Burkart, eliminating
    the level offset in raw posterior-mean RR.
  \item \textbf{PAF contributions are summed signed}, allowing
    sub-TMREL-but-near-optimal exposures to cancel rather than zero.
  \item \textbf{TMREL is year-specific} with nearest-year fill (1990 for years
    $\leq 2000$, 2010 for 2000--2014, 2020 for 2015+) rather than
    averaged across years.
  \item \textbf{Burden $= \text{deaths} \times \text{PAF}$.} The SEV
    multiplicative term is dropped; SEV remains as a diagnostic output but
    does not enter burden.
  \item \textbf{SEV is computed once per year} (no daily summation $\times$ cap
    bug), with the conventional definition.
  \item \textbf{National life table only.} YLLs use the national life table
    broadcast across all subnational units. The age 0 $\to$ 1 remap and
    the \texttt{ex - 2.5} / \texttt{ex = 10} truncations are dropped;
    the full life-expectancy column from the life table is used directly.
\end{enumerate}

\subsection{Depto-day attribution becomes optional}

The depto-day branch in \texttt{05\_compute\_pafs.R} (Section
\ref{subsec:depto-day}) is gated on \texttt{COLOMBIA\_VERIFICATION \&\&\ !USE\_DRAWS \&\&\ <daily mortality file present>}. In production:

\begin{itemize}[itemsep=2pt]
  \item For locations \textbf{with} daily mortality: this branch is currently only enabled in
    verification mode; whether to enable for production is open
    (open decision, not yet made). The production default is annual
    attribution: $\text{burden} = \text{deaths\_annual} \times \text{PAF\_annual}$
    at full \texttt{(year, subloc, age, sex, cause)} granularity.
  \item For locations \textbf{without} daily mortality (most countries):
    annual attribution is the only option.
\end{itemize}

\subsection{Subnational dimension is preserved}

The subnational refactor (PAFs at depto, mortality at depto-age-sex,
merge at depto-age-sex-cause) is part of the production pipeline, not a
verification-only change. For non-Colombia locations, \texttt{subloc\_id}
falls back to a single value equal to \texttt{LOCATION\_ID} if the input
data has no subnational column.

\subsection{Uncertainty propagation}

Verification runs use summary mode (\texttt{USE\_DRAWS = FALSE}) for speed:
mean / lower / upper RR, mean TMREL. Production should run with
\texttt{USE\_DRAWS = TRUE} to propagate the full 1000 ERF draws $\times$ 100
TMREL draws (recycled 10$\times$ to match) through the entire pipeline,
producing 95\% uncertainty intervals on PAFs and burdens.

\subsection{Inputs that change at the global scale}

\begin{itemize}[itemsep=2pt]
  \item \textbf{Mortality:} GBD 2023 cause-specific death estimates (parsed
    via \texttt{util\_parse\_gbd\_mortality.R}) replace DANE vital
    registration. These are annual, not daily.
  \item \textbf{Life tables:} UN World Population Prospects (parsed via
    \texttt{util\_parse\_un\_lifetables.R}) replace DANE life tables.
    Annual, single-region (national) by default.
  \item \textbf{Temperature:} ERA5 reanalysis (historical), CMIP6 climate
    projections (forecast). Daily pixel-level format unchanged.
  \item \textbf{Population:} WorldPop gridded, year-varying, 1km. Same
    format as Colombia.
  \item \textbf{TMREL:} IHME files for 1{,}034 GBD locations, 1990 / 2010
    / 2020 with nearest-year fill.
\end{itemize}

\section{Recommendations and next steps}

\paragraph{Short term (validation closure).}
\begin{enumerate}
  \item Test H2 (median vs.\ mean ERF draws) --- one-line change, immediate
    rerun, cheap.
  \item Test H5 (road-injury age restriction) --- inspect Samuel's script,
    add filter if found, rerun.
  \item Decide whether $-11.76\%$ on a single cause-specific YLL is a
    blocker, given that totals, top-3 rankings, sex split, and other
    per-cause numbers all pass.
\end{enumerate}

\paragraph{Medium term (Samuel asks).}
\begin{enumerate}
  \item Request the pre-imputation mortality file or the imputation
    script (H1).
  \item Request his TMREL source / derivation script (H3).
  \item Confirm whether his \texttt{c\_muerte} field has garbage-code
    redistribution applied and with what coefficients (H4).
\end{enumerate}

\paragraph{Before any production / global run.} Verify
\texttt{COLOMBIA\_VERIFICATION = FALSE} in \texttt{config.R}. The safety
guard catches mis-set flags only when \texttt{LOCATION\_ID != 125}, so
nothing prevents the flag from staying \texttt{TRUE} on a re-run of
Colombia where it would silently apply the wrong methodology to a
production output.

\bigskip
\hrule
\smallskip

\textit{Source artifacts:}
\texttt{output/results/comparison\_to\_samuel\_125.csv} (numerical
comparison),
\texttt{output/results/burden\_125.rds} (full-granularity attributable
burden),
\texttt{output/results/pafs\_daily\_125.rds} (daily PAFs),
\texttt{global-scripts/} (pipeline source on branch
\texttt{verification-depto-day}).

\end{document}
